\chapter{Coordinate systems}\label{ch:crs}
\begin{goals}
\textbf{In this chapter} you give a project its coordinate system, open vector files and point clouds that
were saved in another one, and read what micro did to them and how accurately.
\end{goals}

Survey pickups arrive in SIRGAS 2000, an old lease map in SAD69, a GPS track from the field in WGS 84
longitude and latitude, and the block model in the mine's UTM zone. micro keeps one coordinate system
for the project and moves vectors and points into it as they open, saying each time how they were
moved and how far the result can be trusted.

\section{The project's coordinate system}\label{sec:crs-project}

A project has one \term{coordinate reference system} (CRS). The chip at the right of the status bar,
beside \ui{Sealed}, shows it — \code{EPSG:31983}, or \code{no CRS} — and a click on it opens the selector
(\cref{sec:crs-selector}). The same command is \menu{File>Coordinate system…} and, in the palette,
\key{⌘K} → \code{coordinate system}.

The system is set in one of four ways:
\begin{description}
\item[from a file] the first vector file or point cloud you open with a stated \emph{projected} system (a
  \file{.prj}, a GeoJSON \code{"crs"}, an SRID, a LAS coordinate-system record) gives the project its
  system, when the project has none yet;
\item[suggested] the selector proposes systems whose area of use holds your data;
\item[picked] in the selector, by name or EPSG code;
\item[typed] pasted into the selector: an EPSG code, a proj4 string (\code{+proj=tmerc …}) or a WKT
  definition — what a \file{.prj} holds.
\end{description}

A project's system must be \emph{projected}: micro works in metres, so a geographic system (degrees) is
refused with that reason. The system is saved in \file{project.json} and set again before any layer
reopens. The layer properties' \ui{source} tab and \menu{Data>File info} show it, beside what the file
itself said.

\section{The selector}\label{sec:crs-selector}

\shot{crs-picker}{The coordinate-system selector: SIRGAS 2000 / UTM zone 23S found by name, its area of use (amber) on the world map beside the data's extent (violet), the details below.}

The selector carries the \emph{whole} EPSG table: every projected and geographic system that has a proj4
definition — 6,983 of them, with their names and areas of use. It is part of micro.html and is unpacked
the first time the selector opens; nothing is fetched.

\begin{description}
\item[search] words of the name in any order (\code{sirgas utm 23s}) or a code (\code{31983});
\item[filters] \ui{all}, \ui{projected}, \ui{geographic}, \ui{recent} (the systems you have used), and
  \ui{fits my data} — only systems whose area of use holds the data's extent;
\item[suggested] a row of systems for where the data sits: modern datums, UTM zones and metres first,
  then the tightest area of use;
\item[the map] the world (Natural Earth, 1:110 million) with the selected system's area of use in amber
  and your data's extent in violet;
\item[details] name, code, kind, units, datum, the area of use in words and degrees, the stated accuracy
  of its datum shift to WGS 84, the proj4 definition — and whether your data lies inside its area;
\item[paste] a definition read by \ui{Read}: an ESRI \file{.prj} is matched to its EPSG code by name
  when it carries none (\code{SAD\_1969\_UTM\_Zone\_23S} is EPSG:29193).
\end{description}

Double-click a row, or press \key{Enter} in the search box, to use it.

\begin{tryit}
Find a system by name and by code.
\begin{enumerate}
\item Click the CRS chip in the status bar.
\item Type \code{sirgas utm 23s}: the list narrows to SIRGAS 2000 / UTM zone 23S.
\item Select it. Its area of use — Brazil between 48°W and 42°W — shows in amber on the map.
\item Type \code{31983} instead: the same system, found by its code.
\end{enumerate}
\end{tryit}

\section{Reprojection on read}\label{sec:crs-reproject}

\term{Reprojection} happens when a file opens, and only for the kinds of data where moving every
coordinate is the right answer:

\begin{itemize}
\item strings and polygons in every format micro reads (\cref{ch:polygons}) — CSV string files, WKT,
  WKT-CSV, shapefiles, GeoJSON, KML and KMZ, DXF as features — and points files read from them;
\item point clouds: a LAS file whose coordinate-system record (an OGC WKT, or GeoTIFF keys) names a
  system other than the project's.
\end{itemize}

Block models, grids and meshes are \emph{not} reprojected: moving each block of a model would break its
boxes. They are read in their own coordinates.

What counts as the file's system: its \file{.prj} or SRID or \code{"crs"}; for GeoJSON without a
\code{"crs"} and for KML, WGS 84 longitude and latitude (both formats say so); for a file that says
nothing but whose coordinates all fit within ±180 and ±90 over less than 20°, WGS 84 too. A file that
says nothing and is in metres is taken to be in the project's system already.

The move is horizontal. Each vertex's z is carried unchanged: elevations in mine data are above sea
level, and the change in ellipsoidal height a datum shift implies is not a change of elevation.

\begin{caution}
A longitude/latitude file opened while the project has no system is refused, with an offer to set one
(the selector opens with suggestions for where the file sits). micro never places degrees on a metre
grid.
\end{caution}

\section{How accurate the move is}\label{sec:crs-accuracy}

Every reprojected layer records the path its coordinates took, in the layer's \ui{source} tab, in
\menu{Data>File info}, in its lineage (\code{reproject}) and in the open dialog's line for the file:

\begin{lstlisting}
SAD69 → SIRGAS 2000 by 3-parameter Helmert
  (-67.35, +3.88, -38.22 m, IBGE R.PR 1/2005) — ~1–5 m
\end{lstlisting}

The record carries one of five \term{accuracy classes}:

\begin{table}[htbp]\small
\begin{tabularx}{\linewidth}{@{}lL@{}}\toprule
class & meaning \\ \midrule
\code{exact} & the same datum: projection maths only, millimetres (SIRGAS 2000 degrees → SIRGAS 2000 / UTM) \\
\code{sub-metre} & WGS 84, SIRGAS 2000, ETRS89 and NAD83-class frames taken as identical — they agree to centimetres; no shift is applied \\
\code{metre} & a Helmert datum shift, three or seven parameters, with its source and stated accuracy (SAD69, Córrego Alegre, older national datums) \\
\code{grid} & a datum shift by an NTv2 grid — none is bundled \\
\code{unknown} & the shift needs a grid that is not loaded, or a system states no datum: no shift was applied, and the error may be tens to hundreds of metres \\ \bottomrule
\end{tabularx}
\caption{Accuracy classes of a reprojection.}\label{tab:crs-accuracy}
\end{table}

\begin{why}
Inside Brazil, the national standard for SAD69 and Córrego Alegre is IBGE's: the parameters of
Resolução R.PR 1/2005 (SAD69: ΔX −67.35, ΔY +3.88, ΔZ −38.22 m; Córrego Alegre: ΔX −206.05,
ΔY +168.28, ΔZ −3.82 m). micro uses them for SAD69 data that sits in Brazil, in place of EPSG's
continent-wide mean (−57, +1, −41 m, stated as ~19 m), and for Córrego Alegre in place of a grid it
does not carry. Outside Brazil, SAD69 takes EPSG's default — the record says which.
\end{why}

\section{Exporting in longitude and latitude}\label{sec:crs-export}

A strings and polygons layer exports, besides its own system's formats, two in WGS 84 degrees
(\menu{Export features…}): \ui{KML}, for Google Earth and a GPS, every attribute as ExtendedData; and
\ui{GeoJSON in lon/lat}, the RFC 7946 shape, with no \code{"crs"} member. Both are reprojected from the
layer's system on the way out, so a layer needs one: the project's, or the one its file stated.

A shapefile from a layer that was reprojected on read is written without a \file{.prj}: the original
one names the old system. The status line names the project's system to set in the reading software.

\section{Captions}\label{sec:crs-captions}

Captions and title blocks print the project's system with \code{\{crs\}} —
\code{SIRGAS 2000 / UTM zone 23S · EPSG:31983} — or its halves, \code{\{crs:name\}} and
\code{\{crs:code\}} (\cref{sec:compo-fields}).
